\chapter{The Code Plan: What We Test, How, and What We Have Tested}
\label{ch:codeplan}

Chapter~\ref{ch:selection} chose the methods; this chapter is the engineering contract that keeps them
honest. It answers three questions in order. What kinds of tests does the code base run, and what does each
kind actually certify (Section~\ref{sec:cp-harness})? What has already been tested on the synthetic worlds,
and what did we learn --- including from the tests the data generator itself failed
(Section~\ref{sec:cp-log})? What remains to be run before the calibration campaign is complete, stated as
pass/fail specifications rather than intentions (Section~\ref{sec:cp-forward})? Section~\ref{sec:cp-ci}
closes with the continuous-integration gates and the definition of done. Where a number already appears in a
table of Chapter~\ref{ch:selection} we cite it rather than repeat it; this chapter adds the numbers that
belong to the engineering log rather than to the decision framework.

\section{Testing philosophy and harness}
\label{sec:cp-harness}

A calibration code base fails in two distinct ways: the code computes the wrong number, or the code computes
the right number for a procedure that cannot recover the truth. Conventional software tests catch only the
first. We therefore run four layers of tests, each certifying something the previous layer cannot.

\paragraph{Layer 1: unit tests on hand-computable micro-examples.} Every nontrivial data transformation has
a companion example small enough to verify by hand. Three anchor cases: (i) a three-firm choice set with
known funding histories, where the pools, the excluded most-recently-funded firm, the newcomer row, and the
winner labels produced by \texttt{build\_choice\_sets} are asserted against a worked-out table; (ii) a
two-bucket compensator for an exponential-kernel intensity, checked against the closed-form integral to
machine precision; (iii) deduplication on a constructed deal table containing exact duplicates, 2-day
near-duplicates, and legitimate consecutive-week repeats, asserting which rows survive. These tests certify
arithmetic and indexing, nothing more --- but when a recovery test fails, they are what lets us exonerate the
plumbing in minutes.

\paragraph{Layer 2: equivalence tests between fast and reference fitters.} The package ships slow,
transparent loop fitters and vectorized fast fitters (\texttt{fit\_sector\_glm\_fast},
\texttt{fit\_choice\_fast}). The fast versions exploit structure: the sector GLM is separable per receiving
sector, so twelve small concave problems replace one large one; the ranker is a padded-tensor batch over
events and candidates. Speed is worthless if the optimum moves, so on subsamples both implementations must
agree at the optimum --- same loss to $10^{-6}$, parameters to matching tolerance --- before a fast fitter may
be used in any experiment. The check is rerun whenever either implementation changes.

\paragraph{Layer 3: statistical recovery tests.} These are the scientific tests. We fit on synthetic data
whose ground truth we wrote down, and a run passes when the truth lies within 2 standard errors of the
estimate across 5 seeds. A recovery failure with passing layers 1--2 is a finding about the \emph{method}
--- identifiability, confounding, misspecification --- and several such findings became sections of this
manual. Recovery tests are the only layer that can certify the phrase ``the pipeline works.''

\paragraph{Layer 4: invariant and leakage audits.} Two standing assertions run over the whole feature
stack. The point-in-time assertion recomputes every feature using only data truncated at time $t$ and
requires the result to be bit-identical to the production value: any divergence means some path read the
future. The oracle-isolation assertion requires that no fitting path opens \texttt{ground\_truth.npz} or
\texttt{private\_deal\_map.csv}; oracle files are for evaluation only, and a fit that touches them is void
regardless of its metrics. Both audits are structural: they do not check that a number is right, they check
that it was legal to compute.

\begin{table}[htbp]
\centering
\small
\begin{tabular}{lp{5.6cm}l}
\toprule
Layer & What it certifies & Cadence \\
\midrule
Unit (micro-examples) & arithmetic and indexing of each transformation & every commit \\
Equivalence (fast vs.\ loop) & the vectorized optimum is the reference optimum & every commit \\
Recovery (synthetic truth) & the procedure can find truth within 2 SE, 5 seeds & nightly, small configs \\
Invariant / leakage & features are point-in-time; oracle files untouched & blocking, every commit \\
\bottomrule
\end{tabular}
\caption{The four test layers and what a green result actually means.}
\label{tab:cp-layers}
\end{table}

Three reproducibility rules bind everything above. First, every number reported in this manual is
regenerable from one seeded entry point --- for the worlds themselves,
\begin{verbatim}
PYTHONPATH=. python -m synthetic.generate --regime A --seed 0
\end{verbatim}
and analogously one script per experiment table with a \texttt{--seed} argument. Second, temporal splits
only, everywhere: train weeks $[0,416)$ against test weeks $[416,520)$ on the weekly track, an 80\% day
quantile on the choice track, and no random split anywhere in the stack, including model selection. Third,
runtimes are logged alongside results, because a fitter that silently becomes ten times slower is a
regression even when its numbers do not move.

\section{What we have tested so far}
\label{sec:cp-log}

Everything in this section is measured, seed 0, on the two synthetic worlds of Chapter~\ref{ch:problem}:
regime A, the well-specified weekly two-stage world (50{,}036 deals), and regime B, the deliberately
misspecified daily block Hawkes world (49{,}788 deals), both spanning 520 weeks, 12 sectors, and 8{,}000
firms, and both passed through the full observation-noise stack. The headline recovery tables live in
Chapter~\ref{ch:selection}; here we record the log entries around them --- what broke, what was measured on
the way, and what each episode taught the production design.

\subsection{Tests the world itself failed}
\label{sec:cp-generator}

The generator was not born stable. We present its construction-time failures as what they were: recovery
tests run against the simulator, which the simulator failed until fixed. Each failure is a preview of a
production incident, because every mechanism involved --- log-link feedback, cumulative covariates,
deduplication --- exists in the real pipeline too.

First, metastability. With log-link raw-count feedback at an effective spectral radius of 0.35 ---
comfortably subcritical by linear-Hawkes intuition --- the world was metastable: a single ordinary burst
tipped the roughly 50k-event trajectory into a supercritical excursion that finished at 169{,}436 events.
The shipped generator runs at radius 0.15 with a self-consistent intercept calibration, and the episode is
analyzed as a stability result in Chapter~\ref{ch:selection}. Second, rich-get-richer collapse. With
unbounded cumulative features (raised-to-date, employees) feeding positive selection weights and no exit
mechanism, deal flow concentrated absurdly: one firm collected 422 deals. The fix was structural, not a
tuning knob --- size-damped employee growth capped at 2{,}000 and stage-graduation exits that retire
persistent winners --- after which the maximum per-firm count fell to 32 and the 99th percentile to 17.
Third, deduplication overreach. A naive 7-day same-company window, meant to absorb the 1.2\% injected
duplicates under 6-day date jitter, destroyed roughly 38k legitimate consecutive-week deals --- most of the
repeat-funding signal the model exists to capture. The shipped window is 3 days, validated on the layer-1
constructed table.

\begin{table}[htbp]
\centering
\footnotesize
\begin{tabular}{p{3.2cm}p{4.1cm}p{4.0cm}p{2.9cm}}
\toprule
Failure & Symptom & Fix & Post-fix check \\
\midrule
Log-link metastability & radius 0.35 world tipped from $\sim$50k to 169{,}436 events &
radius 0.15; self-consistent intercept calibration & event counts stable across seeds \\
Rich-get-richer collapse & one firm with 422 deals & size-damped employee growth (cap 2{,}000);
stage-graduation exits & max 32; p99 = 17 \\
Dedup overreach & 7-day window destroyed $\sim$38k legitimate repeat deals & 3-day window &
constructed-table unit test \\
\bottomrule
\end{tabular}
\caption{Generator failures during construction, treated as recovery tests the world initially failed.}
\label{tab:cp-generator}
\end{table}

\begin{remark}
\label{rem:cp-generator}
These are not embarrassments to hide but the cheapest incidents we will ever have. Each cost hours on
synthetic data; the identical mechanisms in production would cost quarters of misallocated attention before
anyone noticed. The production lessons are collected in Table~\ref{tab:cp-lessons}.
\end{remark}

\subsection{Stage-1 recovery: the weekly sector GLM}
\label{sec:cp-stage1}

Training on weeks $[0,416)$, the weekly sector GLM recovers regime A --- the world it is specified for ---
up to observation noise: covariate $\beta$ correlation with truth 0.675, roughly 25\% attenuation from the
errors-in-variables induced by estimated sizes and jittered dates; intercept bias $-0.09$; excitation-entry
correlation 0.27. The one systematic failure is causal, not numerical: the fitted effective spectral radius
is 0.355 against a true 0.172, because the latent serially correlated risk-appetite factor is absorbed as
spurious excitation. On regime B the same tools face a daily block Hawkes world and degrade as they should:
$\beta$ correlation 0.37, fitted radius 0.57 (largely a daily-to-weekly aggregation artifact), and a
held-out Pearson dispersion of 2.72 that correctly flags the misspecification. The full table is in
Chapter~\ref{ch:selection}; the log entry here is the engineering one: the fast fitter reproduces the loop
fitter's optimum (layer 2) and completes in 0.2 seconds where the loop fitter takes minutes, which is what
makes the 5-seed, multi-variant forward plan of Section~\ref{sec:cp-forward} affordable at all.

\subsection{Choice-model recovery: the funded-pool conditional logit}
\label{sec:cp-choice}

The stage-2 model is the funded-pool conditional logit on truthful risk sets with a 64-candidate sampled
softmax \cite{mcfadden1974}, split at the 80\% day quantile. On regime A it reaches a held-out NLL per event
of 3.346 against 4.011 for random choice and a top-5 accuracy of 0.446; on regime B, 3.796 against 4.172
and top-5 of 0.316. The newcomer option wins roughly 27\% of events in both worlds. The exclusion rule
dropped 35 immediate-repeat events in A and 2 in B. Construction of the choice tensors takes 6.4 seconds
and the fit about 12 seconds. The recovered regime-A weights are worth a table, because the pattern of what
recovers and what transmutes is a finding in itself.

\begin{table}[htbp]
\centering
\small
\begin{tabular}{lcc}
\toprule
Feature & Fitted (regime A) & Truth \\
\midrule
Raised-to-date & 0.336 & 0.30 \\
Employees & 0.253 & 0.30 \\
Age & 0.054 & $-0.15$ \\
Stage & 0.132 & 0.15 \\
Log-gap (time since last funding) & $+1.146$ & (26-week cooldown $\approx -2.7$) \\
Newcomer ASC & 6.856 & (prices the entrant pool) \\
\bottomrule
\end{tabular}
\caption{Choice-model weight recovery, regime A, seed 0. Regime B analogues: log-gap $+0.437$, newcomer
ASC 4.691.}
\label{tab:cp-weights}
\end{table}

Raised, employees, and stage recover to within attenuation. The age sign is missed because the true
mechanism there is recency, and the fitted positive log-gap weight is the model's translation of the true
cooldown: a 26-week suppression indicator of about $-2.7$ reappears as ``longer gap, higher score,'' which
is the same economics in a different parametrization. In regime B the softer log-gap weight of $+0.437$
matches the smooth 30-day-half-life inhibition of that world. Whether the log form is the \emph{right}
carrier of this structure --- as opposed to binned dummies or an explicit exponential decay --- is exactly
experiment E2 of the forward plan.

\subsection{Risk-set contamination}
\label{sec:cp-riskset}

One experiment was promoted from finding to assumption. With a naive observed-active mask --- a firm counted
as at risk only when the vendor shows it active --- the ranker does not merely lose likelihood (held-out NLL
3.886 versus 3.558 for the oracle mask); it \emph{inverts} the economics, fitting a recency weight of
$+1.94$ where the truth is 0, because firms enter the observed mask by transacting. The oracle mask recovers
all weights. The full table and the warning it earned are in Chapter~\ref{ch:selection}. The consequence for
the code plan: truthful risk sets are now a stated design assumption, and the contaminated-mask fit is kept
as a standing audit --- if a production risk set is ever inferred rather than given, this is the first
experiment to rerun.

\subsection{Operator smoke test}
\label{sec:cp-operator}

The physics-informed neural operator machinery for the compensator equations runs end to end: data
generation, training, and inference are wired and reproducible. A deliberately cheap 600-epoch demonstration
reaches 8.2\% relative $L^2$ error --- machinery confirmed, accuracy not yet claimed. Under the full training
protocol the operator reaches sub-2\% for 3--8 sector configurations at roughly 3 ms per solve, which is the
regime that matters for the many-repeated-solves role it is assigned in experiment E4.

\subsection{From findings to production rules}
\label{sec:cp-lessons}

\begin{table}[htbp]
\centering
\footnotesize
\begin{tabular}{p{5.6cm}p{7.6cm}}
\toprule
Finding (measured) & Production lesson \\
\midrule
Log-link raw-count feedback metastable at radius 0.35 (50k $\to$ 169{,}436 events) &
guard simulation well below the linear threshold for log-link feedback; prefer bounded transforms of
counts \\
One firm collected 422 deals under unbounded cumulative features with no exits &
cap or damp cumulative covariates; model exits explicitly in any forward-simulation loop \\
7-day dedup destroyed $\sim$38k legitimate repeat deals &
dedup windows are modelling decisions; validate against constructed cases before counting anything \\
Fitted radius 0.355 vs.\ true 0.172 under a latent risk-appetite factor &
never read fitted excitation causally without factor proxies and a sensitivity band \\
Naive risk-set mask inverted the recency weight ($+1.94$ vs.\ true 0) &
treat risk-set truthfulness as an audited input, not a fittable nuisance \\
\bottomrule
\end{tabular}
\caption{The five log entries of this section, restated as production rules.}
\label{tab:cp-lessons}
\end{table}

\section{The forward test plan}
\label{sec:cp-forward}

What follows is the committed test plan: each entry states its purpose, its procedure in two or three
lines, and a pass criterion frozen \emph{before} the run, so that results can fail. Unless stated
otherwise, every experiment runs on both regimes at 5 seeds with error bars, under the splits and audit
rules of Section~\ref{sec:cp-harness}. All entries are currently \emph{planned};
Table~\ref{tab:cp-forward} summarizes.

\paragraph{E1 --- When/where layer: weekly GLM versus event-time Hawkes.}
\emph{Purpose:} quantify what day-resolution timestamps buy over weekly buckets, and how badly latent
confounding contaminates fitted excitation. \emph{Procedure:} fit the weekly GLM and a multivariate
event-time sector Hawkes (exponential kernel, log-linear macro baseline) on the same streams; rerun the
event-time fit with dates bucketed at 1d, 7d, and 28d, reporting bias and SE of the excitation scale
$\kappa$ and decay $\theta$ at each width; separately, regenerate regime A with
\texttt{latent\_strength} $\in \{0, 0.15, 0.3\}$ at fixed true radius and plot fitted effective radius
against latent strength, with and without mitigations (lagged aggregate count covariate, off-diagonal
penalties). \emph{Pass:} event-time fit beats the weekly GLM on held-out NLL on regime B; $\theta$
recovered within 2 SE at day resolution on B while the $\kappa$--$\theta$ ridge visibly reopens by 28d; the
confounding curve is monotone and yields a stated bucket-width recommendation for real data.

\paragraph{E2 --- Which-firm layer: choice-model sub-experiments.}
\emph{Purpose:} settle the four open design choices inside the conditional logit before it is frozen for
production. \emph{Procedure and pass criteria, per sub-experiment:}
\begin{enumerate}
\item \emph{Gap functional form.} Fit log-gap versus binned gap dummies ($<$1m, 1--3m, 3--6m, 6--12m,
$>$12m) versus exponential decay $\exp(-\text{days}/\tau)$ with $\tau \in \{30, 90, 180\}$ days; select by
held-out NLL and plot each implied hazard-versus-gap curve against truth. Pass: the binned fit recovers the
true 26-week step in regime A and the smooth exponential decay in regime B.
\item \emph{Newcomer designs N1--N4.} Compare the bare ASC against the context-covariate ASC (sector heat
and macro features), the representative-entrant construction (plug-in and integrated Gaussian variants),
the nested two-step, and the immigrant/offspring split that aligns first financings with the Hawkes
background stream. Pass: the recommended design calibrates quarterly newcomer share within 2 percentage
points, with incumbent weights materially unchanged.
\item \emph{Sampled-softmax stability.} Refit at \texttt{max\_candidates} $\in \{32, 64, 128\}$ and with
the full denominator on a 10\% subsample. Pass: weights stable across candidate counts, consistent with the
subsampled-denominator consistency result.
\item \emph{Exclusion-rule sensitivity.} Vary the exclude-most-recent rule (sector-scoped, global, none);
count dropped repeats and measure weight shifts. Pass: the effect is pinned with a number (expected
negligible at 2--35 events), and the rule is stated in the production spec either way.
\end{enumerate}

\paragraph{E3 --- Single-layer block Hawkes versus the two-stage factorization.}
\emph{Purpose:} the central architecture decision --- given exact dates and truthful risk sets, what does
the two-stage factorization still buy, and what does it cost in joint-intensity coherence?
\emph{Procedure:} fit the full firm-level event-time block Hawkes (log-link baseline, self-inhibition,
size-normalized sector field) on day-resolution timestamps; evaluate it and the two-stage pipeline on
identical held-out events for joint NLL, next-deal sector prediction, and funded-firm ranking conditional
on the event. \emph{Pass:} on regime B the single-layer fit recovers the self-inhibition ordering with rank
correlation at least 0.9; the head-to-head produces a written recommendation with numbers on all three
metrics, whichever way it falls.

\paragraph{E4 --- Operator benchmark.}
\emph{Purpose:} establish where the neural-operator surrogate pays for itself now that exact event-time MLE
removes its calibration role. \emph{Procedure:} train under the full protocol; report accuracy versus
sector count $M$ and the wall-clock crossover against the exact solvers under repeated forward solves;
benchmark differentiable-solver inverse calibration against the MLE for bias, variance, and wall-clock on
identical data. \emph{Pass:} sub-2\% relative $L^2$ reproduced at full protocol; the crossover point (in
solves and in $M$) is documented; the inverse-calibration table states plainly whether MLE wins at this
scale.

\paragraph{E5 --- Walk-forward backtest with covariate ablation.}
\emph{Purpose:} the end-to-end rehearsal of the production loop under strictly forward information.
\emph{Procedure:} freeze the best when/where and which-firm models; walk forward weekly over the final 20\%
of days, emitting sector intensities, ranked candidate lists, and newcomer probabilities; score top-1/5/10
and MRR against oracle and random, PIT count calibration with the quasi-Poisson dispersion correction, and
lead time (weeks a funded firm spends in the top decile before its deal); rerun with macro covariates
removed and on the \texttt{covariate\_scale} 1.5 world. \emph{Pass:} the pipeline beats random throughout
and degrades gracefully toward oracle; the covariate ablation turns ``covariates help'' into a signed
number in both directions.

\paragraph{B1 --- XGBoost ranking benchmark.}
\emph{Purpose:} a model-free discriminative ceiling for stage 2, per the standing benchmark of
Chapter~\ref{ch:selection}. \emph{Procedure:} train boosted rankers on the identical choice sets and splits
in two variants --- feature parity (exactly the six logit features, isolating functional form) and kitchen
sink (all point-in-time-legal raw fields, with snapshot columns explicitly forbidden); score with per-event
renormalization; modest hyperparameter search with a time-ordered validation tail. \emph{Pass:} the
comparison against top-5 of 0.446 (A) and 0.316 (B) is reported with the frozen decision rule: a material
win means the logit is missing features or nonlinearity, and the remedy is feature work; a narrow result
keeps the interpretable logit.

\paragraph{B2 --- Non-homogeneous Poisson null.}
\emph{Purpose:} show, not assume, that self-excitation earns its complexity at this data scale.
\emph{Procedure:} fit the nested ladder --- covariates-only inhomogeneous Poisson (simulable by thinning
\cite{lewis1979}), covariates plus excitation, full event-time timing --- plus a seasonal-naive baseline;
decompose held-out improvement across the rungs with residual diagnostics per rung \cite{ogata1988}.
\emph{Pass:} the decomposition table exists with error bars over 5 seeds; if the excitation rung
contributes approximately nothing, the recommendation to drop it is written into the decision tables of
Chapter~\ref{ch:selection}.

\paragraph{Causal-pipeline validation.}
\emph{Purpose:} test the causal stack of Chapter~\ref{ch:causal} against known truth, per the protocol
sketched in Chapter~\ref{ch:selection}. \emph{Procedure:} generate a treatment with known selection on
observables and on the latent quality mark; estimate naive, backdoor-adjusted, and AIPW effects; attach
cluster-bootstrap confidence intervals and E-values. \emph{Pass:} the naive estimate is biased in the
selection direction, adjustment removes the observable component, AIPW matches backdoor with better
nuisance robustness, and the residual gap to truth is consistent with the computed E-value --- each asserted
in code, so a failure is a red test, not a discussion.

\begin{table}[htbp]
\centering
\footnotesize
\begin{tabular}{lp{7.9cm}l}
\toprule
Test & Pass criterion (frozen) & Status \\
\midrule
E1 bucketing & $\theta$ within 2 SE at day resolution on B; ridge visible by 28d & planned \\
E1 confounding curve & monotone fitted-radius curve over latent strength; mitigations quantified & planned \\
E2.1 gap form & binned fit recovers 26-week step (A) and smooth decay (B) & planned \\
E2.2 newcomer N1--N4 & quarterly newcomer share within 2pp; incumbent weights stable & planned \\
E2.3 sampled softmax & weights stable across 32/64/128/full-on-subsample & planned \\
E2.4 exclusion rule & sensitivity pinned with a number & planned \\
E3 head-to-head & rank-corr $\ge 0.9$ on B; three-metric comparison written up & planned \\
E4 operator & sub-2\% rel-$L^2$; crossover and inverse-calibration tables & planned \\
E5 backtest & beats random throughout; ablation deltas reported & planned \\
B1 XGBoost & decision rule applied to top-5 0.446 (A) / 0.316 (B) & planned \\
B2 Poisson null & improvement decomposition with error bars & planned \\
Causal validation & all four coded assertions pass against truth & planned \\
\bottomrule
\end{tabular}
\caption{Forward test plan at a glance. All rows: both regimes where applicable, 5 seeds, temporal splits,
audits of Section~\ref{sec:cp-harness} in force.}
\label{tab:cp-forward}
\end{table}

\section{Continuous integration and definition of done}
\label{sec:cp-ci}

The four layers of Section~\ref{sec:cp-harness} map onto three CI gates. On every commit, the unit and
equivalence suites run and must pass; they are fast because the micro-examples are tiny and the equivalence
checks run on subsamples. Nightly, the recovery suite runs on reduced configurations --- fewer weeks, fewer
seeds --- as an early-warning system; a nightly recovery regression blocks the next release, not the next
commit. The leakage audit is a blocking check at both cadences: a commit that makes any feature
non-point-in-time, or that lets a fitting path touch \texttt{ground\_truth.npz} or
\texttt{private\_deal\_map.csv}, does not merge, whatever its metrics say.

The artifact policy is equally plain. Results tables are versioned CSVs, written by the experiment entry
points and diffed like code, so a silently shifted number is a visible change in review. Figures are
regenerated from those CSVs by scripts and are never hand-edited; a figure that cannot be regenerated is
deleted rather than trusted.

The campaign is done when three conditions hold. Every row of Table~\ref{tab:cp-forward} has moved from
\emph{planned} to \emph{pass} or \emph{fail} with a number attached --- not a narrative, a number. The
benchmarks B1 and B2 are reported in the decision tables of Chapter~\ref{ch:selection}, so that the
method choices there cite measured evidence rather than expectation. And every failed pass criterion is
either fixed and rerun, or documented as a limitation alongside those of Chapter~\ref{ch:selection} ---
because a plan whose failures disappear quietly was never a test plan at all.
